\section{Discussion}

\subsection{What the Model Can and Cannot Claim}

The results support a narrow but important conclusion. The recovered fragility classification can be reconstructed exactly from raw survey variables, and the full LASSO-logit model can replicate the screening rule with very high discrimination. However, this is not the same as independently predicting firm distress. Because several predictors are also components of the constructed outcome, the full model primarily validates the mechanics of the index.

The leakage-reduced benchmark is therefore the more substantively meaningful result. Its ROC-AUC of 0.770 and PR-AUC of 0.428 indicate moderate independent predictive signal, not a deployment-ready risk score. This distinction changes the paper's contribution. The manuscript should be read as an audit of a fragility-screening index and as a prototype for transparent predictive screening, rather than as evidence that the model predicts observed closure, bankruptcy, or survival duration. More generally, the paper demonstrates a workflow for auditing constructed policy-screening outcomes in low-data development settings, where index labels, variable coding, and predictive performance can otherwise produce misleading policy claims.

\subsection{Construct Validity and the Fragility Paradox}

Construct validity remains the central limitation. The official variable labels show that the operative rule uses innovation, website status, and employment-size information, while earlier drafts describe some of these variables as revenue volatility and digital-platform difficulty. This ambiguity affects theory, interpretation, and policy translation. In the current version, official labels take precedence. Consequently, the outcome is interpreted as a constructed screening classification that combines modernization and size-related information.

The most theoretically interesting pattern is the fragility paradox of modernization. Variables often treated as signals of capability, such as process innovation and website status, are central to the recovered fragility classification. One interpretation is that modernization may be reactive: firms innovate or establish digital presence because they are under pressure. A second interpretation is growth-stage vulnerability: modernization may bring coordination costs, liquidity strain, and managerial demands. A third interpretation is institutional incompleteness: digital or innovation activity may not translate into resilience when finance, infrastructure, skills, and market systems are uneven. A fourth interpretation is methodological: the paradox may partly reflect how the screening rule was constructed.

These mechanisms cannot be distinguished with the current cross-sectional data. The paper therefore treats them as theory-guided hypotheses for future research rather than as confirmed explanations. A stronger design would link the screening index to later outcomes, such as firm exit, employment contraction, revenue decline, or recovery after disruption.

\subsection{External Validity and Survey Design}

The sample consists of 358 formal firms from the Rwanda 2023 Enterprise Survey. It does not cover the full Rwandan enterprise ecosystem. Informal enterprises, microenterprises, rural firms below the survey eligibility threshold, and firms outside the WBES frame may face different forms of fragility. Predictive performance in the formal-firm sample should therefore not be generalized to all Rwandan SMEs.

Survey design also matters. This paper reports sample-level predictive diagnostics, while the companion weighted-analysis paper estimates formal-firm fragility prevalence under WBES weights. Before national policy claims are made, future work should combine predictive diagnostics with survey-weighted prevalence, design-based uncertainty intervals, and subgroup checks by region, sector, firm size, ownership, and gender where available.

\subsection{Remaining Empirical Work}

The revised analysis adds calibration, PR-AUC, threshold sensitivity, bootstrap confidence intervals, subgroup diagnostics, and a complete-case audit. These additions improve transparency, but they do not eliminate the need for external validation. The calibration evidence is especially important for deployment: the leakage-reduced model may rank firms moderately well, but its predicted probabilities should not be interpreted as calibrated estimates of true fragility risk without recalibration and external validation. Future work should report calibration intercepts and slopes, compare class-weighted and unweighted specifications, and test Platt scaling or isotonic calibration on validation folds before presenting predicted probabilities as risk probabilities. Future work should also validate the index against observed firm outcomes, test the alternative finance-digital Omega index, compare LASSO with other interpretable statistical-learning models only when those models are reproducibly generated, and audit fairness across policy-relevant groups. The most valuable extension is linkage to Rwanda's Establishment Census and Integrated Business Enterprise Survey, which can clarify what the WBES formal-firm frame omits about informal and microenterprise fragility \citep{nisr2025ibes, nisr2024establishment}.
